% ===================================================================
%  TABEL Nr. 7 — Die dorp in die winter. — Die plaaswerf in die lente.
%  Meme regime que les precedents. Premiere de la DEUXIEME SERIE.
%
%  LE MALE, LA FEMELLE ET LE PETIT SONT TROIS MOTS, ET L'IDO LES
%  FABRIQUE : kavalo / kavalino / kavalyuno. L'afrikaans a des mots
%  pleins — perd, merrie, vulletjie ; bul, koei, kalf ; ram, ooi,
%  lam ; bok, sy, boklam — et chacun porte son renvoi.
%
%  ET « pastoro » EST LE BERGER, NON LE PASTEUR : « herder », avec sa
%  houlette et son chien. Le faux ami est d'autant plus facile en
%  afrikaans que « pastoor » y existe et designe le cure — celui du
%  tableau 3.
%
%  LE DIMINUTIF PORTE ICI UNE DISTINCTION, non un sentiment :
%  « vulletjie » est le poulain et « kuikentjie » le poussin ; ce ne
%  sont pas de petits chevaux ni de petites poules.
% ===================================================================

%%K t07-noto-1 noto
\VUnotes{143.2mm}{%
(*) Sien vir die besonderhede van die maaltyd, tabelle 6 en 13.%
}

%%K t07-apar-1 apar
\VUsaut{4.0mm}
\VUcentre{13.2pt}{90}{TWEEDE REEKS}
\VUsaut{3.0mm}
\VUfilet{16mm}
\VUsaut{5.6mm}
\VUtitre{15.0pt}{60}{TABEL Nr. 7}
\VUsaut{3.0mm}

%%K t07-tit-1 sub
\VUcentre{12.0pt}{0}{\textit{Eerste tafereel.}}
\VUsaut{3.0mm}

%%K t07-tit-2 sub
\VUpk{10.2pt}{40}{Die dorp in die winter.}
\VUsaut{4.0mm}

%%K t07-c1-01-1 p
1. --- Tydens die nuwejaarsvakansie het ek my vader vergesel na Nuwestad, 'n dorpie waar
hy 'n paar handelsake moes afhandel.

%%K t07-c1-02-1 p
2. --- Ons het die \VUgras{poskoets} \textsuperscript{(11)} gebruik wat tussen die stad
en daardie dorpie ry; maar omdat dit baie koud was, was daardie reis in 'n min gerieflike
voertuig nie baie aangenaam nie.

%%K t07-c1-03-1 p
3. --- Daarom het ons 'n ware plesier gevoel toe ons die \VUgras{koetsier} sy wa op die
plein sien stilhou het, voor die \VUgras{herberg} \textsuperscript{(2)} van die
\textit{Wit Beer}. Die \VUgras{herbergier} \textsuperscript{(4)} het op die drumpel van
sy huis op ons gewag terwyl hy kalm sy \VUgras{pyp} gerook het.

%%K t07-c1-03-2 p
Hy het ons genooi om in te kom en ek het my nie laat vra om my voor die takkiesvuur te
gaan sit wat in die vuurherd geknetter het nie. Toe het ek erg gespot met die skerp
\VUgras{noordewind} wat die ou \VUgras{uithangbord} \textsuperscript{(3)} laat kraak het
en \VUgras{die rook} \textsuperscript{(1)} wat uit die skoorsteen kom in swart warrelinge
weggedra het.

%%K t07-c1-04-1 p
4. --- Dit was die uur van die middagete. Sonder om langer te wag, het ons die eenvoudige
maar smaaklike maaltyd wat vir ons bedien is goed geëet (*).

%%K t07-tit-3 sub
\VUpk{10.2pt}{40}{'n Paar besoeke.}
\VUsaut{3.0mm}

%%K t07-c1-05-1 p
5. --- Na die ete het ek my vader vergesel, wat 'n versekeringsagent is, na sy kliënte.
Eers het ons die \VUgras{smid} \textsuperscript{(5)} besoek, wat naby die herberg woon.
Hy was besig om 'n perd te beslaan. Ek het die krag van sy maat bewonder wat die
\VUgras{hoef} van die perd stewig op sy leer voorskoot gehou het, terwyl die smid die
\VUgras{hoefyster} \textsuperscript{(6)} met sterk hamerhoue vasgemaak het.

%%K t07-c1-06-1 p
6. --- Daarna het ons na die \VUgras{skoenmaker} \textsuperscript{(7)} van die dorp
gegaan, die ou Jakobus. Hoewel hy 'n pragtige \VUgras{stewel} as uithangbord het, is hy
tevrede om ou gate-skoene van nuwe sole te voorsien.

%%K t07-c1-06-2 p
Die ou man het lag-lag vir ons gesê: « In die stad was ek 'n "skoenmaker" en ek het geen
kliënte gehad nie. Hier is ek 'n "skoenlapper" en die werk is baie. »

%%K t07-c1-07-1 p
7. --- Hy het met ons gepraat oor sy buurman die \VUgras{barbier} \textsuperscript{(8)},
oor wie die kliënte nie tevrede is nie. Sy \VUgras{skeermesse} is altyd geskaar en sy
\VUgras{skêre} sny nooit goed nie.

Maar omdat hy die enigste barbier van die dorp is, is 'n mens natuurlik gedwing om hom
jou te laat skeer en jou hare te laat sny. Buitendien is hy 'n gierige en onvriendelike
man.

%%K t07-c1-08-1 p
8. --- Gelukkig lyk die ander \VUgras{bure} nie soos hy nie. Byvoorbeeld die
\VUgras{wamaker} \textsuperscript{(9)} is 'n goedhartige man, werksaam en hulpvaardig.
Dit is 'n plesier om 'n wa deur hom te laat regmaak. Sy werk is altyd afgewerk.

%%K t07-c1-09-1 p
9. --- Ou Jakobus het ook nie oor die \VUgras{saalmaker} \textsuperscript{(10)} gekla
nie, wat hy soms help om die \VUgras{tuie} te stik wanneer die werk dring.

%%K t07-c1-09-2 p
My vader het hom oor sy gesin gevra: « Almal gaan dit goed. My kleindogter is in die
\VUgras{skool} \textsuperscript{(12)}. Haar \VUgras{onderwyseres} \textsuperscript{(13)}
is baie tevrede oor haar, sy is 'n goeie \VUgras{leerling} \textsuperscript{(14)}. Sy is
nie soos my stoute seun nie; hy dink net aan speel. Die \VUgras{onderwyser}
\textsuperscript{(15)} kry dit nie reg om hom iets te leer nie. »

%%K t07-tit-4 sub
\VUsaut{3.0mm}
\VUpk{10.2pt}{40}{Dorpsgesels.}
\VUsaut{3.0mm}

%%K t07-c1-10-1 p
10. --- Die ou man het ons daarna die nuus van die dorp vertel. Hulle gaan 'n nuwe skool
bou, want dié wat in die \VUgras{gemeentehuis} ingerig is, het te klein geword.
Buitendien is dit maklik, want dit sal genoeg wees om 'n deel van die \VUgras{meulenaar}
se tuin te koop. Dit sal baie voordelig wees vir die arme man. Weens die koue kan die
\VUgras{meulstene} van die \VUgras{meul} \textsuperscript{(16)} nie meer die koring maal
nie, want die \VUgras{wiel} \textsuperscript{(17)} is deur die \VUgras{ys}
\textsuperscript{(25)} van die \VUgras{stroompie} \textsuperscript{(23)} vasgevries.

%%K t07-c1-11-1 p
11. --- « Wat geboue betref, het u ons nuwe \VUgras{kerk} \textsuperscript{(18)} gesien?
Dit is baie skilderagtig onder sy sneeumantel. Die \VUgras{kraaie} \textsuperscript{(19)},
wat hulle ou kloktoring nie meer herken nie, sit treurig op die groot boom van die
\VUgras{kerkhof} \textsuperscript{(20)}. »

%%K t07-c1-12-1 p
12. --- Daardie gedagte aan die kerkhof het die skoenmaker aan die mense herinner wat my
vader geken het en wat verlede jaar gesterf het. « Hoeveel \VUgras{grafte}
\textsuperscript{(21)}, my goeie meneer, is by die ander gevoeg, onder die
\VUgras{sipres} \textsuperscript{(22)}! Die dood het ons ou notaris weggemaai. Al die
dorpenaars was by sy \VUgras{begrafnis}. »

%%K t07-c1-13-1 p
13. --- « Hier is sy opvolger wat op die stroompie skaats. Hy het pas gekom om die riem
van sy \VUgras{skaats} \textsuperscript{(26)} deur my te laat regmaak, wat gebreek het. »

%%K t07-c1-14-1 p
14. --- « Hier is ook op die oewer van die stroompie die nuwe \VUgras{veldwagter}
\textsuperscript{(27)}. Hy is 'n oud-gendarm wat die deugniete van die dorp skrik aanjaag.
Hulle vlug sodra hulle sy \VUgras{kamaste} \textsuperscript{(29)} en sy \VUgras{kepie}
\textsuperscript{(28)} sien. »

%%K t07-c1-15-1 p
15. --- « Ha! hier is die ou Josef wat 'n \VUgras{karretjie met takkebosse}
\textsuperscript{(30)} vir die bakker bring. --- Dit is waar, het my vader gesê, ek
herken sy span \VUgras{muile} \textsuperscript{(31)}. Ek moet juis met hom praat, ek sal
die geleentheid gebruik. Tot siens, vader Jakobus. »

%%K t07-c1-16-1 p
16. --- Net toe ons uitgaan, het die kinders van die dorp die skool verlaat en
hardloop-hardloop na die \VUgras{glybaan} \textsuperscript{(33)} gestorm met die gevaar
om te val
en 'n lid te breek in die \VUgras{val} \textsuperscript{(34)}. Een van hulle het sy
\VUgras{slee} \textsuperscript{(32)} by die wamaker gaan haal, een van sy maats daarin
laat sit, en hardloop-hardloop vertrek.

%%K t07-c1-17-1 p
17. --- Ander het na 'n \VUgras{sneeuhoop} \textsuperscript{(37)} naby die
\VUgras{brug} \textsuperscript{(40)} gestorm en die \VUgras{sneeuman}
\textsuperscript{(39)} begin bombardeer wat hulle die vorige dag opgerig het en waarop
hulle 'n hoed, 'n pyp en 'n skooltas oor die skouer gesit het.

%%K t07-c1-18-1 p
18. --- Hulle gee min om vir die vriesende wind en die koue. Die meeste het nie eers 'n
\VUgras{serp} \textsuperscript{(35)} nie en, om hulle weer warm te maak na die geveg met
\VUgras{sneeuballe} \textsuperscript{(38)}, is 'n handvol baie warm \VUgras{kastaiings}
\textsuperscript{(42)} vir hulle genoeg, gekoop by die ou \VUgras{verkoopster}
Jakobina, wat van koue bewe onder haar te dun \VUgras{tjalie} \textsuperscript{(43)}.

%%K t07-tit-5 sub
\VUsaut{3.0mm}
\VUcentre{12.0pt}{0}{\textit{Tweede tafereel.}}
\VUsaut{3.0mm}

%%K t07-tit-6 sub
\VUpk{10.2pt}{40}{Die plaaswerf in die lente.}
\VUsaut{4.0mm}

%%K t07-c2-01-1 p
1. --- Ondanks al die plesiere van die winter groet ons met diepe vreugde die terugkoms
van die \VUgras{lente}.

%%K t07-c2-01-2 p
Hoe 'n vrolike prentjie lyk 'n plaaswerf, saans, in die maand Mei!

%%K t07-c2-02-1 p
2. --- Aan die horison gaan die \VUgras{son} \textsuperscript{(1)} stadig lê en oorstroom
die \VUgras{platteland} \textsuperscript{(3)} met sy purper \VUgras{strale}
\textsuperscript{(2)}. Die ywerige \VUgras{ploeër} \textsuperscript{(6)} haas hom om sy
\VUgras{land} \textsuperscript{(4)} te ploeg.

%%K t07-c2-02-2 p
Met sy \VUgras{ploeg} \textsuperscript{(5)} wat aan 'n kragtige \VUgras{perd}
\textsuperscript{(7)} ingespan is, trek hy die \VUgras{voor} \textsuperscript{(8)} waarin
die \VUgras{saaier} \textsuperscript{(9)} met vol hand die \VUgras{saad} sal gooi wat die
goue are sal lewer.

%%K t07-c2-03-1 p
3. --- Die \VUgras{maaiers} \textsuperscript{(11)} met hulle sense het die gras van die
weiland gemaai, die hooiers het die \VUgras{hooi} \textsuperscript{(10)} laat droog word
deur dit met \VUgras{vurke} \textsuperscript{(12)} en harke om te keer, en hier kom hulle
terug.

%%K t07-c2-03-2 p
Sommige loop voor die swaar wa wat deur osse getrek word; hulle dra hulle gereedskap op
die skouer. Ander, luier, het hulle gerieflik op die geurige hooi gaan lê.

%%K t07-c2-04-1 p
4. --- Terwyl hulle verbygaan, groet hulle vriendelik die \VUgras{tuinier}
\textsuperscript{(16)} wat in die \VUgras{boord} \textsuperscript{(13)} besig is om die
\VUgras{vrugtebome} te snoei en die \VUgras{peerbome} \textsuperscript{(14)} te ent. Die
\VUgras{pruimbome} \textsuperscript{(15)} begin blom, en in die goed bemeste
\VUgras{groentetuin} \textsuperscript{(17)} staan die groente, kool en slaai reeds goed.

%%K t07-c2-04-2 p
Op die muurtjie sprei 'n mooi \VUgras{pou} \textsuperscript{(18)} sy stert oop om sy
pragtige vere te laat bewonder.

%%K t07-tit-7 sub
\VUpk{10.2pt}{40}{Die plaaswerf.}
\VUsaut{3.0mm}

%%K t07-c2-05-1 p
5. --- Die deure van die \VUgras{stal} \textsuperscript{(19)} is oop en 'n mens sien die
ruif en die \VUgras{strooisel} \textsuperscript{(23)} van die \VUgras{merrie}
\textsuperscript{(21)} wat die \VUgras{staljong} \textsuperscript{(20)} pas uitgespan
het. Die \VUgras{vulletjie} \textsuperscript{(22)}, so komiek met sy lang bene, volg sy
moeder met bokspronge.

%%K t07-c2-06-1 p
6. --- Naby die \VUgras{put} \textsuperscript{(29)} gaan die \VUgras{koeie}
\textsuperscript{(25)} in die hout \VUgras{trog} \textsuperscript{(26)} drink wat vir
hulle as drinkplek dien. Die \VUgras{kalf} \textsuperscript{(24)} gebruik die
geleentheid om te suip, en die \VUgras{bul} \textsuperscript{(27)}, wat die goeie geur
van die voer ruik, bulk vrolik en lig trots sy kop met sy magtige \VUgras{horings}
\textsuperscript{(28)} op.

%%K t07-c2-07-1 p
7. --- Almal werk. Die \VUgras{diensmeisie} \textsuperscript{(32)} hou die \VUgras{tou}
\textsuperscript{(31)} van die put in die hand. Sy skep water met 'n \VUgras{emmer}
\textsuperscript{(30)} om die vee te laat drink. Naby die \VUgras{skuur}
\textsuperscript{(33)} hark 'n man die strooi bymekaar en voor die deur van die
\VUgras{woonhuis} \textsuperscript{(34)} spin 'n ou \VUgras{spinster}
\textsuperscript{(35)} met haar spinnewiel en haar \VUgras{spinrok} die \VUgras{vlas} of
die \VUgras{hennep} soos in die ou tyd.

%%K t07-tit-8 sub
\VUsaut{3.0mm}
\VUpk{10.2pt}{40}{Die boer.}
\VUsaut{3.0mm}

%%K t07-c2-08-1 p
8. --- Die \VUgras{boer} \textsuperscript{(36)} lei al daardie werk en gee opdragte aan
sy knegte. Hy beveel aan om die \VUgras{eg} \textsuperscript{(40)} en die \VUgras{pik}
\textsuperscript{(39)} onder dak te bring wat 'n mens naby die \VUgras{hok}
\textsuperscript{(38)} van die \VUgras{hond} \textsuperscript{(37)} vergeet het.

%%K t07-c2-09-1 p
9. --- Sy geroep laat 'n \VUgras{duif} \textsuperscript{(41)} skrik, wat met volle
vlerke in die \VUgras{duiwehok} \textsuperscript{(42)} invlieg, en die \VUgras{henne}
\textsuperscript{(43)} wat haas om in die \VUgras{hoenderhok} \textsuperscript{(44)}
terug te gaan.

%%K t07-c2-10-1 p
10. --- Voor die \VUgras{skaapkraal} \textsuperscript{(48)}, hier is die ou
\VUgras{herder} \textsuperscript{(45)} wat sy \VUgras{trop} \textsuperscript{(49)}
terugbring. Terwyl hy sy \VUgras{herderstaf} \textsuperscript{(46)} vashou, wat hom nooit
ontbreek nie, net so min as sy verbleikte \VUgras{bloes} \textsuperscript{(47)}, lei hy
na die kraal die \VUgras{ramme} \textsuperscript{(50)}, \VUgras{ooie}
\textsuperscript{(53)}, \VUgras{lammers} \textsuperscript{(54)}, \VUgras{bokke}
\textsuperscript{(51)} en \VUgras{boklammers} \textsuperscript{(52)}. Sy getroue
\VUgras{hond} \textsuperscript{(55)} Argus kom laaste en spoor met sy geblaf die
agterblyers aan.

%%K t07-c2-11-1 p
11. --- Al daardie geraas en beweging steur nie die rus van die goeie \VUgras{melkkoei}
\textsuperscript{(56)} wat haar \VUgras{klokkie} \textsuperscript{(57)} laat klingel nie,
terwyl 'n mens haar voor die deur van die \VUgras{stal} \textsuperscript{(58)} melk.

%%K t07-c2-12-1 p
12. Sy lyk of sy verstaan dat sy haar melk moet gee, sodat die \VUgras{boerin}
\textsuperscript{(60)} \VUgras{botter} in die \VUgras{karring} \textsuperscript{(61)} kan
gereedmaak of die uitstekende \VUgras{kase} \textsuperscript{(64)} wat van die plank drup
wat op die \VUgras{kuip} \textsuperscript{(63)} gesit is.

%%K t07-c2-13-1 p
13. --- Die hele \VUgras{melkery} \textsuperscript{(59)} word baie skoon gehou deur die
\VUgras{plaaskneg} \textsuperscript{(65)} wat pas die \VUgras{melkkanne}
\textsuperscript{(62)} in die groot emmer gewas het wat hy in die hand dra.

%%K t07-tit-9 sub
\VUsaut{3.0mm}
\VUpk{10.2pt}{40}{Die hoenderwerf.}
\VUsaut{3.0mm}

%%K t07-c2-14-1 p
14. --- Met sy dik houtskoene loop hy 'n bietjie lomp, en so jaag hy die
\VUgras{kalkoen} \textsuperscript{(68)}, die \VUgras{eende} \textsuperscript{(66)}, die
\VUgras{eendmannetjies} \textsuperscript{(67)} en die \VUgras{ganse}
\textsuperscript{(69)} weg wat \VUgras{hulle bekke} \textsuperscript{(70)} wyd oopmaak en
onrustig skreeu.

%%K t07-c2-14-2 p
Die \VUgras{haan} \textsuperscript{(71)}, meer strydlustig, rig sy rooi \VUgras{kam}
\textsuperscript{(72)} op, skud die vere van sy \VUgras{stert} \textsuperscript{(73)} en
laat 'n klinkende « kukeleku » hoor.

%%K t07-c2-15-1 p
15. --- 'n \VUgras{Klukhen} \textsuperscript{(74)} krap op die \VUgras{mishoop}
\textsuperscript{(81)} met haar \VUgras{kuikentjies} \textsuperscript{(75)}. Die
\VUgras{varkie} \textsuperscript{(78)} vlug na sy moeder, die enorme \VUgras{sog}
\textsuperscript{(77)} wat ons naby die varkhok sien.

%%K t07-c2-16-1 p
16. --- In hulle hokke eet die \VUgras{konyne} \textsuperscript{(79)} gulsig die fyn
\VUgras{gras} \textsuperscript{(80)} wat die boer se dogter vir hulle bring. Johanna is
baie lief vir haar konyne en, nadat sy elke dag die vars \VUgras{eiers}
\textsuperscript{(76)} in die hoenderhok gaan haal het, kom sy haar maatjies besoek.
\VUsaut{6.0mm}
\VUfilet{22mm}
